\documentclass[../../main.tex]{subfiles}

\begin{document}

\section{两个守恒律的方程组}\label{section:ch11-两个守恒律}

本节更深入分析 $m=2$ 时的初值问题，即一对守恒律：
\begin{align}
\begin{cases}
u_{1,t}+F_1(u_1,u_2)_x=0,\\
u_{2,t}+F_2(u_1,u_2)_x=0,
\end{cases}
\text{在 } \mathbb{R}\times(0,\infty) \text{ 中}
\label{align:ch11-两个守恒律组}
\end{align}
\begin{align*}
u_1=g_1,\;u_2=g_2 \quad \text{在 } \mathbb{R}\times\{t=0\} \text{ 上}.
\end{align*}
这里 $\mathbf{F}=(F_1,F_2)$，$\mathbf{g}=(g_1,g_2)$，$\mathbf{u}=(u_1,u_2)$。

\subsection{Riemann 不变量}\label{subsection:ch11-Riemann不变量}

我们的意图首先是证明通过适当的因变量非线性变换，可将\eqref{align:ch11-两个守恒律组} 化为更简单的形式。想法是找到两个函数 $w_1,w_2:\mathbb{R}^2\to\mathbb{R}$，沿稀疏波曲线 $R_1,R_2$ 具有良好性质：

\begin{definition}\label{definition:ch11-Riemann不变量}
称
\begin{align*}
w_i:\mathbb{R}^2\to\mathbb{R}
\end{align*}
为第 $i$ Riemann 不变量，只要
\begin{align}
Dw_i(z) \text{ 平行于 } \mathbf{l}_j(z) \quad (z\in\mathbb{R}^2,\;i\neq j).
\label{align:ch11-Riemann不变量条件}
\end{align}
\end{definition}

我们稍后将看到条件\eqref{align:ch11-Riemann不变量条件} 多么有用，但先停下来问 Riemann 不变量是否存在。由于现在取 $m=2$，这很容易。事实上，因为 $\mathbf{l}_j(z)\cdot\mathbf{r}_i(z)=0$（$i\neq j$），我们看到\eqref{align:ch11-Riemann不变量条件} 在 $\mathbb{R}^2$ 中等价于
\begin{align}
Dw_i(z)\cdot\mathbf{r}_i(z)=0 \quad (i=1,2,\;z\in\mathbb{R}^2),
\label{align:ch11-Riemann不变量等价}
\end{align}
即
\begin{align}
w_i \text{ 沿稀疏波曲线 } R_i \text{ 为常数} \quad (i=1,2).
\label{align:ch11-Riemann不变量常数}
\end{align}

特别地，任何满足\eqref{align:ch11-Riemann不变量常数} 的光滑函数 $w_i$ 也满足\eqref{align:ch11-Riemann不变量等价}、\eqref{align:ch11-Riemann不变量条件}，因而是第 $i$ Riemann 不变量。

在 $m>2$ 的情形，Riemann 不变量一般不存在。

现在我们可将 $\mathbf{w}=(w_1,w_2)=(w_1(z_1,z_2),w_2(z_1,z_2))$ 视为状态空间 $\mathbb{R}^2$ 上的新坐标，代替 $\mathbf{z}=(z_1,z_2)$。更精确地，定义 $\mathbf{w}:\mathbb{R}^2\to\mathbb{R}^2$ 为
\begin{align}
\mathbf{w}(z)=\mathbf{w}(z_1,z_2)=(w_1(z_1,z_2),w_2(z_1,z_2)).
\label{align:ch11-w映射}
\end{align}
逆映射为 $\mathbf{z}(\mathbf{w})=\mathbf{z}(w_1,w_2)=(z_1(w_1,w_2),z_2(w_1,w_2))$。

现在利用变换\eqref{align:ch11-w映射} 简化两个守恒律组\eqref{align:ch11-两个守恒律组}。假设 $\mathbf{u}=(u_1,u_2)$ 是\eqref{align:ch11-两个守恒律组} 的光滑解。改变因变量，设
\begin{align}
\mathbf{v}(x,t)\triangleq\mathbf{w}(\mathbf{u}(x,t)) \quad (x\in\mathbb{R},\;t>0).
\label{align:ch11-v定义}
\end{align}
$\mathbf{v}=(v_1,v_2)$ 满足什么 PDE 组？

\begin{theorem}[守恒律与 Riemann 不变量。]\label{theorem:ch11-守恒律与Riemann不变量}函数 $v_1,v_2$ 求解方程组
\begin{align}
\begin{cases}
v_{1,t}+\lambda_2(\mathbf{u})v_{1,x}=0,\\
v_{2,t}+\lambda_1(\mathbf{u})v_{2,x}=0,
\end{cases}
\text{在 } \mathbb{R}\times(0,\infty) \text{ 中}.
\label{align:ch11-Riemann不变量方程组}
\end{align}
要点是方程组\eqref{align:ch11-Riemann不变量方程组} 虽非守恒律形式，但在许多方面比\eqref{align:ch11-两个守恒律组} 简单得多。特别注意，$u_1$ 的 PDE 涉及 $u_{2,x}$ 项，而 $v_1$ 的 PDE 不涉及 $v_{2,x}$。类似地，$v_2$ 的 PDE 不涉及 $v_{1,x}$。
\end{theorem}

\begin{proof}
由\eqref{align:ch11-v定义}，对 $i=1,2$，$i\neq j$，
\begin{align*}
v_{i,t}+\lambda_j(\mathbf{u})v_{i,x}=Dw_i(\mathbf{u})\cdot\mathbf{u}_t+\lambda_j(\mathbf{u})Dw_i(\mathbf{u})\cdot\mathbf{u}_x=Dw_i(\mathbf{u})\cdot(-\mathbf{F}(\mathbf{u})_x+\lambda_j(\mathbf{u})\mathbf{u}_x)=Dw_i(\mathbf{u})\cdot(-D\mathbf{F}(\mathbf{u})+\lambda_j(\mathbf{u})I)\mathbf{u}_x=0,
\end{align*}
因为由定义 $Dw_i$ 平行于 $\mathbf{l}_j$。
\end{proof}

\begin{remark}
(i) 我们可通过引入 ODE
\begin{align}
\dot{x}_i(s)=\lambda_j(\mathbf{u}(x_i(s),s)) \quad (s\geqslant 0)
\label{align:ch11-特征ODE}
\end{align}
来解释 PDE 组\eqref{align:ch11-Riemann不变量方程组}，其中 $i=1,2$，$j\neq i$。由\eqref{align:ch11-Riemann不变量方程组}，
\begin{align}
v_i \text{ 沿曲线 } (x_i(s),s) \text{ 为常数} \quad (s\geqslant 0)
\label{align:ch11-v常数}
\end{align}
对 $i=1,2$。

(ii) 回忆 \S11.2 中真非线性条件为
\begin{align}
D\lambda_i(z)\cdot\mathbf{r}_i(z)\neq 0 \quad (z\in\mathbb{R}^2,\;i=1,2).
\label{align:ch11-真非线性条件}
\end{align}
由于我们也可将 $\lambda_i$ 视为 $\mathbf{w}=(w_1,w_2)$ 的函数，可将\eqref{align:ch11-真非线性条件} 改写为
\begin{align}
\frac{\partial\lambda_i}{\partial w_j}\neq 0 \quad (\mathbf{w}\in\mathbb{R}^2,\;i\neq j).
\label{align:ch11-真非线性w}
\end{align}
为看出\eqref{align:ch11-真非线性条件} 等价于\eqref{align:ch11-真非线性w}，注意若\eqref{align:ch11-真非线性w} 失效，则
\begin{align*}
0=\frac{\partial\lambda_i}{\partial w_j}=\sum_{k=1}^2\frac{\partial\lambda_i}{\partial z_k}\frac{\partial z_k}{\partial w_j}.
\end{align*}
但由于 $\sum_{k=1}^2\dfrac{\partial w_i}{\partial z_k}\dfrac{\partial z_k}{\partial w_j}=\delta_{ij}=0$（$i\neq j$），我们看到\eqref{align:ch11-真非线性w} 断言 $D\lambda_i$ 平行于 $Dw_i$。然而 $Dw_i$ 垂直于 $\mathbf{r}_i$，故与\eqref{align:ch11-真非线性条件} 矛盾。因此\eqref{align:ch11-真非线性条件} 蕴含\eqref{align:ch11-真非线性w}，反之亦然。
\end{remark}

\begin{example}[正压可压缩气体动力学。]\label{example:ch11-正压气体动力学}我们通过详细考察可压缩气体动力学 Euler 方程（\S11.1\cref{example:ch11-Euler方程}）在内能 $e$ 为常数的特殊情形来说明上述思想。相关 PDE 为
\begin{align}
\begin{cases}
\rho_t+(\rho v)_x=0 & \text{（质量守恒）},\\
(\rho v)_t+(\rho v^2+p)_x=0 & \text{（动量守恒）},
\end{cases}
\label{align:ch11-正压Euler}
\end{align}
其中假设
\begin{align}
p=p(\rho)
\label{align:ch11-正压状态方程}
\end{align}
对某光滑函数 $p:\mathbb{R}\to\mathbb{R}$。公式\eqref{align:ch11-正压状态方程} 称为\textbf{正压状态方程}。我们假设严格双曲性条件
\begin{align}
p'>0.
\label{align:ch11-正压双曲}
\end{align}
设 $\mathbf{u}=(u_1,u_2)=(\rho,\rho v)$，可将\eqref{align:ch11-正压Euler}、\eqref{align:ch11-正压状态方程} 写为 $\mathbf{u}_t+\mathbf{F}(\mathbf{u})_x=0$，其中
\begin{align*}
\mathbf{F}=(F_1,F_2)=\left(z_2,\frac{(z_2)^2}{z_1}+p(z_1)\right)
\end{align*}
且 $\mathbf{z}=(z_1,z_2)$，前提是 $z_1>0$。则
\begin{align*}
D\mathbf{F}=\begin{pmatrix}0 & 1\\ -\left(\dfrac{z_2}{z_1}\right)^2+p'(z_1) & \dfrac{2z_2}{z_1}\end{pmatrix}.
\end{align*}
因此
\begin{align}
\lambda_1=\frac{z_2}{z_1}-p'(z_1)^{1/2},\quad \lambda_2=\frac{z_2}{z_1}+p'(z_1)^{1/2}.
\label{align:ch11-正压特征值}
\end{align}
用物理记号，
\begin{align}
\lambda_1=v-\sigma,\quad \lambda_2=v+\sigma,
\label{align:ch11-正压物理特征值}
\end{align}
其中声速
\begin{align}
\sigma\triangleq p'(\rho)^{1/2}.
\label{align:ch11-声速}
\end{align}

回忆\eqref{align:ch11-特征ODE}，接下来考虑 ODE
\begin{align}
\dot{x}_1(t)=v(x_1(t),t)+\sigma(x_1(t),t),
\label{align:ch11-特征x1}
\end{align}
\begin{align}
\dot{x}_2(t)=v(x_2(t),t)-\sigma(x_2(t),t),
\label{align:ch11-特征x2}
\end{align}
其中 $\sigma(x,t)\triangleq p'(\rho(x,t))^{1/2}$，$t\geqslant 0$。由\eqref{align:ch11-v常数}，Riemann 不变量 $v_1=w_1(\mathbf{u})$ 沿\eqref{align:ch11-特征x1} 的轨迹为常数，$v_2=w_2(\mathbf{u})$ 沿\eqref{align:ch11-特征x2} 的轨迹为常数。

为直接计算 $w_1$ 和 $w_2$，对系统\eqref{align:ch11-正压Euler} 进行一些计算。首先将\eqref{align:ch11-正压Euler} 化为非散度形式：
\begin{align}
\rho_t+\rho v_x+\rho_x v=0,
\label{align:ch11-质量非散度}
\end{align}
\begin{align}
\rho_t v+\rho v_t+\rho_x v^2+2\rho v v_x+p_x=0.
\label{align:ch11-动量非散度}
\end{align}
将\eqref{align:ch11-质量非散度} 乘以 $\sigma^2=p'(\rho)$ 并回忆\eqref{align:ch11-正压状态方程}，得
\begin{align}
p_t+vp_x+\sigma^2\rho v_x=0.
\label{align:ch11-压强方程}
\end{align}
此外\eqref{align:ch11-质量非散度}、\eqref{align:ch11-动量非散度} 合并得
\begin{align}
\rho v_t+\rho v v_x+p_x=0.
\label{align:ch11-动量简化}
\end{align}

现在操纵\eqref{align:ch11-压强方程}、\eqref{align:ch11-动量简化} 使方向 $\lambda_1,\lambda_2=v\mp\sigma$ 显式出现。为此将\eqref{align:ch11-动量简化} 乘以 $\sigma$，然后加到\eqref{align:ch11-压强方程} 并从\eqref{align:ch11-压强方程} 减去：
\begin{align}
\begin{cases}
p_t+(v+\sigma)p_x+\rho\sigma(v_t+(v+\sigma)v_x)=0,\\
p_t+(v-\sigma)p_x-\rho\sigma(v_t+(v-\sigma)v_x)=0.
\end{cases}
\label{align:ch11-特征组合}
\end{align}
由此推出
\begin{align}
\begin{cases}
\dfrac{\mathrm{d}}{\mathrm{d}t}[p(x_1(t),t)]+\rho(x_1(t),t)\sigma(x_1(t),t)\dfrac{\mathrm{d}}{\mathrm{d}t}[v(x_1(t),t)]=0,\\
\dfrac{\mathrm{d}}{\mathrm{d}t}[p(x_2(t),t)]-\rho(x_2(t),t)\sigma(x_2(t),t)\dfrac{\mathrm{d}}{\mathrm{d}t}[v(x_2(t),t)]=0.
\end{cases}
\label{align:ch11-沿轨迹方程}
\end{align}
由于 $\dfrac{\mathrm{d}p}{\mathrm{d}t}=\sigma^2\dfrac{\mathrm{d}\rho}{\mathrm{d}t}$，我们看到
\begin{align}
\frac{\sigma}{\rho}\frac{\mathrm{d}\rho}{\mathrm{d}t}\pm\frac{\mathrm{d}v}{\mathrm{d}t}=0 \quad \text{沿\eqref{align:ch11-特征x1}、\eqref{align:ch11-特征x2} 的轨迹},
\label{align:ch11-Riemann不变量微分}
\end{align}
前提是 $\rho>0$。

现在将 Riemann 不变量视为 $\rho$ 和 $v$ 的函数。由于 $v_1=w_1(\rho,v)$ 沿由 $x_1(\cdot)$ 确定的曲线为常数，
\begin{align*}
0=\frac{\mathrm{d}}{\mathrm{d}t}[w_1(\rho(x_1(t),t),v(x_1(t),t))]=\frac{\partial w_1}{\partial\rho}\frac{\mathrm{d}}{\mathrm{d}t}[\rho(x_1(t),t)]+\frac{\partial w_1}{\partial v}\frac{\mathrm{d}}{\mathrm{d}t}[v(x_1(t),t)].
\end{align*}
这与\eqref{align:ch11-Riemann不变量微分} 一致，若
\begin{align*}
\frac{\partial w_1}{\partial\rho}=\frac{\sigma(\rho)}{\rho},\quad \frac{\partial w_1}{\partial v}=1.
\end{align*}
类似地推出
\begin{align*}
\frac{\partial w_2}{\partial\rho}=\frac{\sigma(\rho)}{\rho},\quad \frac{\partial w_2}{\partial v}=-1.
\end{align*}
积分后，我们得出 Riemann 不变量（相差加性常数）为
\begin{align*}
w_1=\int_1^\rho\frac{\sigma(s)}{s}\,\mathrm{d}s+v,\quad w_2=\int_1^\rho\frac{\sigma(s)}{s}\,\mathrm{d}s-v.
\end{align*}
验证 $w_1,w_2$ 作为 $\mathbf{z}=(z_1,z_2)$ 的函数满足 Riemann 不变量定义留作练习。
\end{example}

\subsection{光滑解的不存在性}\label{subsection:ch11-光滑解不存在}

现在说明 Riemann 不变量的用处，建立光滑解不存在的以下判据：

\begin{theorem}[Riemann 不变量与爆破。]\label{theorem:ch11-Riemann不变量与爆破}假设 $\mathbf{g}$ 光滑、具有紧支集。进一步假设真非线性条件
\begin{align}
\frac{\partial\lambda_i}{\partial w_j}>0 \quad \text{在 } \mathbb{R}^2 \text{ 中} \quad (i=1,2,\;i\neq j)
\label{align:ch11-真非线性正}
\end{align}
成立。则初值问题\eqref{align:ch11-两个守恒律组} 不可能存在对所有 $t\geqslant 0$ 都存在的光滑解 $\mathbf{u}$，若
\begin{align}
\text{要么 } v_{1,x}<0 \text{ 要么 } v_{2,x}<0 \quad \text{在 } \mathbb{R}\times\{t=0\} \text{ 上某处}.
\label{align:ch11-爆破条件}
\end{align}
\end{theorem}

\begin{proof}
1. 暂时假设 $\mathbf{u}$ 是\eqref{align:ch11-两个守恒律组} 的光滑解。写
\begin{align}
a\triangleq v_{1,x},\quad b\triangleq v_{2,x},
\label{align:ch11-a定义}
\end{align}
其中 $\mathbf{v}=\mathbf{w}(\mathbf{u})$，$\mathbf{v}=(v_1,v_2)$ 求解 PDE 组\eqref{align:ch11-Riemann不变量方程组}。对\eqref{align:ch11-Riemann不变量方程组} 的第一个方程关于 $x$ 求导，计算
\begin{align}
a_t+\lambda_2 a_x+\frac{\partial\lambda_2}{\partial w_1}a^2+\frac{\partial\lambda_2}{\partial w_2}ab=0.
\label{align:ch11-a方程}
\end{align}
利用\eqref{align:ch11-Riemann不变量方程组} 的第二个方程，改写为
\begin{align*}
v_{2,t}+\lambda_2 v_{2,x}=(\lambda_2-\lambda_1)b.
\end{align*}
代入\eqref{align:ch11-a方程} 得
\begin{align}
a_t+\lambda_2 a_x+\frac{\partial\lambda_2}{\partial w_1}a^2+\left[\frac{1}{\lambda_2-\lambda_1}\frac{\partial\lambda_2}{\partial w_2}(v_{2,t}+\lambda_2 v_{2,x})\right]a=0.
\label{align:ch11-a方程2}
\end{align}

2. 为积分\eqref{align:ch11-a方程2}，固定 $x_0\in\mathbb{R}$，设
\begin{align}
\xi(t)\triangleq\exp\left(\int_0^t\frac{1}{\lambda_2-\lambda_1}\frac{\partial\lambda_2}{\partial w_2}(v_{2,t}+\lambda_2 v_{2,x})(x_1(s),s)\,\mathrm{d}s\right),
\label{align:ch11-xi定义}
\end{align}
其中
\begin{align*}
\begin{cases}
\dot{x}_1(s)=\lambda_2(\mathbf{u}(x_1(s),s)) & (s\geqslant 0),\\
x_1(0)=x_0.
\end{cases}
\end{align*}
由\eqref{align:ch11-v常数} 的关键观察：$v_1$ 沿曲线 $(x_1(s),s)$ 为常数。写
\begin{align*}
v_1(x_1(s),s)=v_1^0=v_1(x_0,0) \quad (s\geqslant 0).
\end{align*}
因此表达式 $\left(\dfrac{1}{\lambda_2-\lambda_1}\dfrac{\partial\lambda_2}{\partial w_2}\right)$ 作为 $\mathbf{v}=\mathbf{w}(\mathbf{u})$ 的函数，仅依赖于 $v_2$。设
\begin{align*}
\gamma(\mu)\triangleq\int_0^\mu\left(\frac{1}{\lambda_2-\lambda_1}\frac{\partial\lambda_2}{\partial w_2}\right)(v_1^0,v)\,\mathrm{d}v.
\end{align*}
则\eqref{align:ch11-xi定义} 蕴含
\begin{align}
\xi(t)=\exp\left(\int_0^t\frac{\mathrm{d}}{\mathrm{d}s}[\gamma(v_2(x_1(s),s))]\,\mathrm{d}s\right)=\exp(\gamma(v_2(x_1(t),t))-\gamma(v_2(x_0,0))).
\label{align:ch11-xi化简}
\end{align}

3. 将\eqref{align:ch11-a方程2} 变换为
\begin{align*}
\frac{\mathrm{d}}{\mathrm{d}t}((\xi\alpha)^{-1})=-\frac{1}{(\xi\alpha)^2}\frac{\mathrm{d}}{\mathrm{d}t}(\xi\alpha)=\frac{\partial\lambda_2}{\partial w_1}\xi^{-1},
\end{align*}
其中 $\alpha(t)\coloneqq a(x_1(t),t)$，假设 $\alpha\neq 0$。因此
\begin{align*}
(\xi(t)\alpha(t))^{-1}=(\alpha(0))^{-1}+\int_0^t\frac{\partial\lambda_2}{\partial w_1}\xi^{-1}(s)\,\mathrm{d}s.
\end{align*}
此等式重排为
\begin{align}
\alpha(t)=\alpha(0)\xi^{-1}(t)\left(1+\alpha(0)\int_0^t\frac{\partial\lambda_2}{\partial w_1}\xi^{-1}(s)\,\mathrm{d}s\right)^{-1}.
\label{align:ch11-alpha公式}
\end{align}

4. 由 PDE 组\eqref{align:ch11-Riemann不变量方程组}，$\mathbf{v}$ 有界。由\eqref{align:ch11-xi化简}，存在常数 $\theta,\Theta$ 使得 $0<\theta\leqslant\xi(t)\leqslant\Theta$ 对所有 $t>0$。因此由\eqref{align:ch11-真非线性正} 和\eqref{align:ch11-alpha公式}，$\alpha$ 对所有 $t>0$ 有界当且仅当 $\alpha(0)\geqslant 0$，即 $v_{1,x}(x_0,0)\geqslant 0$。

类似计算将 $v_1$ 换为 $v_2$。我们得出：若在 $\mathbb{R}\times\{t=0\}$ 上某处 $v_{1,x}<0$ 或 $v_{2,x}<0$，则不存在\eqref{align:ch11-两个守恒律组} 的对所有 $t\geqslant 0$ 都持续的光滑解。
\end{proof}

\end{document}